% GENERATED FILE. Edit canonical lessons, then run npm run generate:books.
% canonical-source-hash: fnv1a64:4f85eb3545036562
% canonical-lessons: BN-C304-utsob, BN-C304-pujo, BN-C304-protima, BN-C304-pendel, BN-C304-udjapon-kora

\chapter{Festival Time}
\label{ch:bn-festival-time}
\hlchaptermodality{304}

\begin{chapteropening}
\textbf{By the end of this chapter:} \emph{I can talk about a festival, its images and its pandals.}

Complete the last lesson of chapter 304: I can talk about a festival, its images and its pandals.
\end{chapteropening}

\section[\texorpdfstring{\bn{উৎসব} (utsôb)}{উৎসব (utsôb)}]{\bn{উৎসব} (utsôb) — a festival}

\label{lesson:BN-C304-utsob}



\begin{quote}
Before the new one: say the Bengali for in the course of time, then the Bengali for to commemorate.
\end{quote}

\subsection*{You'll want to know: \bn{উৎসব}}
\textbf{\bn{উৎসব}} — \emph{utsôb} — \textquotedblleft{}a festival\textquotedblright{}.

\begin{quote}
\bn{এই} \bn{গ্রামে} \bn{বসন্তে} \bn{একটা} \bn{বড়} \bn{উৎসব} \bn{হয়।}
\end{quote}

\begin{quote}
\emph{ei grāme bôsônte ekṭā bôṛo utsôb hôy.}
\end{quote}

\begin{quote}
There is a big festival in this village in spring.
\end{quote}

\begin{cousinweb}
\textbf{\bn{উৎসব}} (\emph{utsôb}) is Sanskrit \emph{utsava}, first 'an undertaking, a beginning' and then 'a festival' (Monier-Williams).
\end{cousinweb}

\begin{grammarlens}[title={What you've built}]
One more word for telling what happened.
\end{grammarlens}

\subsection*{Your turn}
\begin{itemize}
\raggedright
  \item \emph{Say it:} \emph{utsôb}
  \item \emph{Say it:} \emph{utsôb}, once more
  \item \emph{Say it:} \emph{ei grāme bôsônte ekṭā bôṛo utsôb hôy.} — There is a big festival in this village in spring.
  \item \emph{Recall:} say the Bengali for in the course of time, then the Bengali for to commemorate, then say \emph{utsôb} again
\end{itemize}

\begin{tcolorbox}[breakable,colback=teal!4,colframe=teal!35!black,title={Before you move on}]
What does \bn{উৎসব} mean? (\textquotedblleft{}A festival\textquotedblright{}.) Say it once more.
\end{tcolorbox}
\section[\texorpdfstring{\bn{পুজো} (pujo)}{পুজো (pujo)}]{\bn{পুজো} (pujo) — a puja; a worship festival}

\label{lesson:BN-C304-pujo}



\begin{quote}
Before the new one: say the Bengali for to commemorate, then the Bengali for a festival.
\end{quote}

\subsection*{You'll want to know: \bn{পুজো}}
\textbf{\bn{পুজো}} — \emph{pujo} — \textquotedblleft{}a puja; a worship festival\textquotedblright{}.

\begin{quote}
\bn{পুজোর} \bn{সময়} \bn{সবাই} \bn{নতুন} \bn{জামা} \bn{পরে।}
\end{quote}

\begin{quote}
\emph{pujor sômôy sôbāi nôtun jāmā pôre.}
\end{quote}

\begin{quote}
At puja time everyone wears new clothes.
\end{quote}

\begin{cousinweb}
\textbf{\bn{পুজো}} (\emph{pujo}) is the spoken form of Sanskrit \emph{pūjā}, 'honour, worship' (Monier-Williams), which the written word \bn{পূজা} keeps. In Bengal \bn{পুজো} alone means the great autumn festival.
\end{cousinweb}

\begin{grammarlens}[title={What you've built}]
One more word for telling what happened.
\end{grammarlens}

\subsection*{Your turn}
\begin{itemize}
\raggedright
  \item \emph{Say it:} \emph{pujo}
  \item \emph{Say it:} \emph{pujo}, once more
  \item \emph{Say it:} \emph{pujor sômôy sôbāi nôtun jāmā pôre.} — At puja time everyone wears new clothes.
  \item \emph{Recall:} say the Bengali for to commemorate, then the Bengali for a festival, then say \emph{pujo} again
\end{itemize}

\begin{tcolorbox}[breakable,colback=teal!4,colframe=teal!35!black,title={Before you move on}]
What does \bn{পুজো} mean? (\textquotedblleft{}A puja; a worship festival\textquotedblright{}.) Say it once more.
\end{tcolorbox}
\section[\texorpdfstring{\bn{প্রতিমা} (prôtimā)}{প্রতিমা (prôtimā)}]{\bn{প্রতিমা} (prôtimā) — an image of a deity; an idol}

\label{lesson:BN-C304-protima}



\begin{quote}
Before the new one: say the Bengali for a festival, then the Bengali for a puja.
\end{quote}

\subsection*{You'll want to know: \bn{প্রতিমা}}
\textbf{\bn{প্রতিমা}} — \emph{prôtimā} — \textquotedblleft{}an image of a deity; an idol\textquotedblright{}.

\begin{quote}
\bn{কুমোররা} \bn{মাটি} \bn{দিয়ে} \bn{প্রতিমা} \bn{বানায়।}
\end{quote}

\begin{quote}
\emph{kumorrā māṭi diye prôtimā bānāy.}
\end{quote}

\begin{quote}
Potters make the images out of clay.
\end{quote}

\begin{cousinweb}
\textbf{\bn{প্রতিমা}} (\emph{prôtimā}) is Sanskrit \emph{pratimā}, 'an image, a likeness; a statue' (Monier-Williams), from \emph{prati-√mā}, 'to copy'.
\end{cousinweb}

\begin{grammarlens}[title={What you've built}]
One more word for telling what happened.
\end{grammarlens}

\subsection*{Your turn}
\begin{itemize}
\raggedright
  \item \emph{Say it:} \emph{prôtimā}
  \item \emph{Say it:} \emph{prôtimā}, once more
  \item \emph{Say it:} \emph{kumorrā māṭi diye prôtimā bānāy.} — Potters make the images out of clay.
  \item \emph{Recall:} say the Bengali for a festival, then the Bengali for a puja, then say \emph{prôtimā} again
\end{itemize}

\begin{tcolorbox}[breakable,colback=teal!4,colframe=teal!35!black,title={Before you move on}]
What does \bn{প্রতিমা} mean? (\textquotedblleft{}An image of a deity; an idol\textquotedblright{}.) Say it once more.
\end{tcolorbox}
\section[\texorpdfstring{\bn{প্যান্ডেল} (pænḍel)}{প্যান্ডেল (pænḍel)}]{\bn{প্যান্ডেল} (pænḍel) — a pandal; a festival marquee}

\label{lesson:BN-C304-pendel}



\begin{quote}
Before the new one: say the Bengali for a puja, then the Bengali for an idol.
\end{quote}

\subsection*{You'll want to know: \bn{প্যান্ডেল}}
\textbf{\bn{প্যান্ডেল}} — \emph{pænḍel} — \textquotedblleft{}a pandal; a festival marquee\textquotedblright{}.

\begin{quote}
\bn{পাড়ার} \bn{প্যান্ডেলে} \bn{রাতে} \bn{অনেক} \bn{ভিড়} \bn{হয়।}
\end{quote}

\begin{quote}
\emph{pāṛār pænḍele rāte ônek bhiṛ hôy.}
\end{quote}

\begin{quote}
At night there's a big crowd at the neighbourhood pandal.
\end{quote}

\begin{cousinweb}
\textbf{\bn{প্যান্ডেল}} (\emph{pænḍel}) is Indian English \emph{pandal}, a marquee of bamboo and cloth, a word usually traced to Tamil \emph{pantal}, 'a shed'. Its initial \bn{প্যা} is read \emph{pæ}.
\end{cousinweb}

\begin{grammarlens}[title={What you've built}]
One more word for telling what happened.
\end{grammarlens}

\subsection*{Your turn}
\begin{itemize}
\raggedright
  \item \emph{Say it:} \emph{pænḍel}
  \item \emph{Say it:} \emph{pænḍel}, once more
  \item \emph{Say it:} \emph{pāṛār pænḍele rāte ônek bhiṛ hôy.} — At night there's a big crowd at the neighbourhood pandal.
  \item \emph{Recall:} say the Bengali for a puja, then the Bengali for an idol, then say \emph{pænḍel} again
\end{itemize}

\begin{tcolorbox}[breakable,colback=teal!4,colframe=teal!35!black,title={Before you move on}]
What does \bn{প্যান্ডেল} mean? (\textquotedblleft{}A pandal; a festival marquee\textquotedblright{}.) Say it once more.
\end{tcolorbox}
\section[\texorpdfstring{\bn{উদযাপন} \bn{করা} (udjāpôn kôrā)}{উদযাপন করা (udjāpôn kôrā)}]{\bn{উদযাপন} \bn{করা} (udjāpôn kôrā) — to celebrate}

\label{lesson:BN-C304-udjapon-kora}



\begin{quote}
Before the new one: say the Bengali for an idol, then the Bengali for a pandal.
\end{quote}

\subsection*{You'll want to know: \bn{উদযাপন} \bn{করা}}
\textbf{\bn{উদযাপন} \bn{করা}} — \emph{udjāpôn kôrā} — \textquotedblleft{}to celebrate\textquotedblright{}.

\begin{quote}
\bn{গ্রামের} \bn{মানুষ} \bn{নাচ} \bn{আর} \bn{গান} \bn{দিয়ে} \bn{উৎসব} \bn{উদযাপন} \bn{করল।}
\end{quote}

\begin{quote}
\emph{grāmer mānush nāch ār gān diye utsôb udjāpôn kôrlô.}
\end{quote}

\begin{quote}
The villagers celebrated the festival with dancing and songs.
\end{quote}

Past (he/she …): \bn{তিনি} \bn{জন্মদিন} \bn{উদযাপন} \bn{করলেন}, \emph{tini jônmôdin udjāpôn kôrlen}.

\begin{cousinweb}
\textbf{\bn{উদযাপন} \bn{করা}} (\emph{udjāpôn kôrā}) puts Sanskrit \emph{udyāpana}, 'bringing to a conclusion' (Monier-Williams), the closing rite of a vow, before \bn{করা}; Bengali uses it for celebrating.
\end{cousinweb}

\begin{grammarlens}[title={What you've built}]
One more word for telling what happened.
\end{grammarlens}

\subsection*{Your turn}
\begin{itemize}
\raggedright
  \item \emph{Say it:} \emph{udjāpôn kôrā}
  \item \emph{Say it:} \emph{udjāpôn kôrā}, once more
  \item \emph{Say it:} \emph{grāmer mānush nāch ār gān diye utsôb udjāpôn kôrlô.} — The villagers celebrated the festival with dancing and songs.
  \item \emph{Recall:} say the Bengali for an idol, then the Bengali for a pandal, then say \emph{udjāpôn kôrā} again
\end{itemize}

\begin{tcolorbox}[breakable,colback=teal!4,colframe=teal!35!black,title={Before you move on}]
What does \bn{উদযাপন} \bn{করা} mean? (\textquotedblleft{}To celebrate\textquotedblright{}.) Say it once more.
\end{tcolorbox}
